\documentclass[11pt,a4paper]{article}

% ──────────────────────────────────────────────────────────────────────────
% AXON-RFC-006 Appendix C — OpenAI Compatibility as Agent Dispatch View.
%                            Capability-URI API keys.
% Status: Draft v0.1.
%
%   RFC-006-C is the second reference realisation of the paradigm
%   RFC-006-B v0.6 made explicit: external protocols are
%   transport views over the invocation graph, not new
%   ontologies.
%
%   RFC-006-B (Pages): HTTP file fetches → ability invocations.
%   RFC-006-C (this):  OpenAI-compatible LLM endpoints → agent
%                      dispatch streams. The chat-base output of
%                      any dispatched agent (Claude Code, Codex,
%                      future LLM drivers) is, after a small
%                      stream filter, the wire shape of an OpenAI
%                      `/v1/chat/completions` response. EasyNet
%                      therefore exposes every `<agent>.chat`
%                      ability as an OpenAI-compatible model
%                      without writing a new chat surface — the
%                      adapter is a view, not a wrap.
%
%   API keys are minted by keyring abilities and live as
%   `resource/api_key.<id>` URAs. Persisting an API key is a
%   capability-style affair: holding the token = capability;
%   revoking the resource = capability gone. There is no OAuth
%   layer, no token exchange, no scopes-as-bag-of-strings.
%   Receipts attribute every OpenAI request to the api_key URA
%   and chain forward through the agent's dispatch.
%
%   Three normative invariants centre v0.1:
%     INV-1  Adapter Purity     — `01HUB.openai.chat_completions`
%                                 is a pure transport view; no
%                                 cached state; no canonical-
%                                 receipt mutation outside the
%                                 single auth-trail entry.
%     INV-2  Capability-URI Key — API keys are minted by keyring
%                                 abilities, addressed as
%                                 `resource/api_key.<id>`,
%                                 and self-authoritative
%                                 (presence = grant; absence =
%                                 deny).
%     INV-3  Filter Determinism — the stream-json → SSE filter
%                                 is deterministic over the agent
%                                 dispatch sequence.
%     INV-4  Auth Receipt Trail — every request mints a canonical
%                                 receipt over the API-key
%                                 resource; subsequent dispatch
%                                 receipts cite it as
%                                 `causal_context.Scalar`.
%
% Build:    xelatex AXON-RFC-006-C-openai-compat.tex
% Companion (paradigm sibling):   AXON-RFC-006-B-easynet-webapp.tex (v0.6)
% Companion (axiom):              AXON-RFC-006-stateful-easynet.tex
% Companion (URA + ontology):     AXON-RFC-001-plan-v4.1.5.md
% Companion (ontology theorem):   EasyNet-Axon/document/concepts/ONTOLOGY_AGENT_ABILITY.md
% Companion (keyring):            EasyNet-Cli/docs/rfc/AXON-RFC-002-stage-1-3-edit-plan.md
% ──────────────────────────────────────────────────────────────────────────

\usepackage[margin=1in]{geometry}
\usepackage{amsmath,amssymb,amsthm}
\usepackage{xcolor}
\usepackage[colorlinks=true,linkcolor=blue!60!black,citecolor=blue!60!black,urlcolor=blue!60!black]{hyperref}
\usepackage{booktabs}
\usepackage{longtable}
\usepackage{array}
\usepackage{tabularx}
\usepackage{xltabular}
\newcolumntype{L}[1]{>{\raggedright\arraybackslash}p{#1}}
\usepackage{listings}
\usepackage{enumitem}
\usepackage{tcolorbox}
\usepackage{titlesec}
\usepackage{fancyhdr}
\usepackage{parskip}
\usepackage{xspace}

\usepackage{fontspec}
\usepackage{xeCJK}
\setCJKmainfont{PingFang SC}[
    BoldFont={PingFang SC Semibold},
    ItalicFont={PingFang SC Light}
]

\pagestyle{fancy}
\fancyhf{}
\fancyhead[L]{\small AXON-RFC-006 Appendix C v0.1}
\fancyhead[R]{\small OpenAI compat as dispatch view}
\fancyfoot[C]{\thepage}
\renewcommand{\headrulewidth}{0.4pt}

\titleformat{\section}{\Large\bfseries\color{black!85}}{\thesection}{0.7em}{}
\titleformat{\subsection}{\large\bfseries\color{black!75}}{\thesubsection}{0.6em}{}
\titleformat{\subsubsection}{\normalsize\bfseries\color{black!70}}{\thesubsubsection}{0.5em}{}

\definecolor{normativecolor}{RGB}{180,30,30}
\definecolor{derivedcolor}{RGB}{30,90,180}
\definecolor{deferredcolor}{RGB}{120,90,30}
\definecolor{codebg}{RGB}{245,245,248}
\definecolor{codeborder}{RGB}{210,210,220}

\newcommand{\normsv}{\textcolor{normativecolor}{\textbf{[normative]}}\xspace}
\newcommand{\derived}{\textcolor{derivedcolor}{\textbf{[derived]}}\xspace}
\newcommand{\deferred}{\textcolor{deferredcolor}{\textbf{[deferred]}}\xspace}

\lstdefinestyle{ealstyle}{
    backgroundcolor=\color{codebg},
    basicstyle=\ttfamily\small,
    frame=single,
    rulecolor=\color{codeborder},
    breaklines=true,
    showstringspaces=false,
    keywordstyle=\color{blue!70!black}\bfseries,
    commentstyle=\color{green!40!black}\itshape,
    stringstyle=\color{red!60!black},
    morekeywords={ability,resource,user,agent,hub,kind,realm,owner,uri,subject,caller,callee,receipt,call,let,mission,with,timeout,on,fn,struct,impl,pub,async,await,assume,axiom,invariant,capability},
    columns=fullflexible,
    keepspaces=true,
    aboveskip=0.6em,
    belowskip=0.6em,
}
\lstset{style=ealstyle}

\newtheorem{theorem}{Theorem}
\newtheorem{lemma}[theorem]{Lemma}
\newtheorem{corollary}[theorem]{Corollary}
\theoremstyle{definition}
\newtheorem{definition}{Definition}

\newtcolorbox{invariant}[1][Invariant]{
    colback=red!4!white,
    colframe=red!50!black,
    fonttitle=\bfseries,
    title={Invariant --- #1},
    sharp corners=south,
    boxrule=0.6pt,
}
\newtcolorbox{notebox}[1][]{
    colback=blue!4!white,
    colframe=blue!50!black,
    fonttitle=\bfseries,
    title={Note},
    sharp corners=south,
    boxrule=0.6pt,
    #1
}
\newtcolorbox{warningbox}[1][]{
    colback=orange!6!white,
    colframe=orange!70!black,
    fonttitle=\bfseries,
    title={Boundary},
    sharp corners=south,
    boxrule=0.6pt,
    #1
}
\newtcolorbox{decisionbox}[1][]{
    colback=purple!4!white,
    colframe=purple!60!black,
    fonttitle=\bfseries,
    title={Open decision},
    sharp corners=south,
    boxrule=0.6pt,
    #1
}
\newtcolorbox{assumptionbox}[1]{
    colback=green!4!white,
    colframe=green!50!black,
    fonttitle=\bfseries,
    title={Assumption --- #1},
    sharp corners=south,
    boxrule=0.6pt,
}

\title{
    \vspace{-2em}
    \textbf{\Large AXON-RFC-006 Appendix C (v0.1):}\\[0.2em]
    \textbf{\Large OpenAI Compatibility as Agent Dispatch View,}\\[0.1em]
    \textbf{\Large Capability-URI API Keys}\\[0.4em]
    \large LLM endpoints are not a new ontology.\\
    \large Every \texttt{<agent>.chat} ability is an OpenAI-shape\\
    \large \texttt{/v1/chat/completions} model after one stream filter.
}
\author{
    Silan Hu \\ \small National University of Singapore
}
\date{Draft v0.1, \today}

\begin{document}
\maketitle

\begin{abstract}
\noindent
RFC-006-B v0.6 established the discipline that external protocols
enter EasyNet as transport views over the invocation graph rather
than as new ontologies. Pages was the first reference instance:
HTTP file fetches become deterministic projections of the
\texttt{<user>.<project>.page.fetch} ability. This appendix
specifies the second reference instance: an OpenAI-compatible
\texttt{/v1/chat/completions} endpoint exposed by the Hub on top
of any agent's chat-base ability output. The construction is not
a wrap. The dispatched agent --- Claude Code, Codex, or any
future driver --- already emits a streaming sequence of typed
frames; those frames, after a small deterministic filter, are
exactly the wire shape of an OpenAI streaming completion. The
adapter is a view. Authentication is not OAuth: API keys are
minted by keyring abilities and addressed as
\texttt{resource/api\_key.<id>} URAs; presence of the token at
the boundary is presence of a capability, and revocation is a
single keyring call. Four normative invariants centre v0.1:
Adapter Purity (the OpenAI surface mutates nothing it would not
have mutated through the regular dispatch path), Capability-URI
Key (API keys are URAs, not OAuth tokens), Filter Determinism
(the stream-json $\to$ SSE projection is deterministic over the
dispatch sequence), and Auth Receipt Trail (every request mints
one canonical receipt against the API-key resource, all
subsequent dispatch receipts cite it). The MVP is two endpoints
on the Hub --- \texttt{/v1/chat/completions} and
\texttt{/v1/models} --- backed by a fixed-shape stream filter
implemented in the Hub agent's \texttt{openai.chat\_completions}
ability. Cursor, Continue, openai-python, and any other
OpenAI-shape client work against EasyNet without code changes.
\end{abstract}

\tableofcontents
\vspace{1em}
\hrule
\vspace{1em}

% ──────────────────────────────────────────────────────────────────────────
\section{Position in the RFC-006 series}
\label{sec:position}

\derived{} RFC-006-B v0.6 made one paradigm explicit:

\begin{quote}
External protocols are transport views over the invocation
graph, not new ontologies.
\end{quote}

It then realised the paradigm against one external protocol ---
HTTP file fetch --- and called the realisation \emph{Pages}.
This appendix realises the same paradigm against a second
external protocol --- the OpenAI \texttt{/v1/chat/completions}
streaming wire format --- and calls the realisation
\emph{OpenAI compat}. The structure of the appendix mirrors
RFC-006-B section by section, so a reader can read the two
side by side and verify that the same disciplines apply.

\begin{center}
\renewcommand{\arraystretch}{1.25}
\begin{tabularx}{\linewidth}{l X X}
\toprule
\textbf{Aspect} & \textbf{RFC-006-B (Pages)} & \textbf{RFC-006-C (this)} \\
\midrule
External protocol           & HTTP / GET / static files                          & OpenAI \texttt{/v1/chat/completions} (streaming) \\
Underlying ability          & \texttt{<user>.<project>.page.fetch}              & \texttt{<agent>.chat} (or any chat-base ability) \\
Hub adapter ability         & \texttt{01HUB.pages.serve}                         & \texttt{01HUB.openai.chat\_completions} \\
URL form                    & \texttt{<project>.<user>.pages.<realm>/<path>}    & \texttt{<realm>/v1/chat/completions} \\
Is the adapter pure?        & Yes (RFC-006-B INV-1)                              & Yes (this RFC INV-1) \\
Identity addressing         & \texttt{resource/<user>.<project>/<path>}          & \texttt{resource/api\_key.<id>} \\
Determinism guarantee       & Same id $\to$ same bytes (RFC-006-B INV-3)        & Same dispatch frames $\to$ same SSE (this RFC INV-3) \\
What ships in the MVP       & two files (hello-world.html + style.css)           & one streaming completion against \texttt{web-builder.chat} \\
\bottomrule
\end{tabularx}
\end{center}

The two appendices share the same vocabulary, the same kind of
invariant numbers, the same closing claim. RFC-006-D, RFC-006-E,
RFC-006-F (file storage, conversation snapshots, mission outputs)
will follow the same structure.

% ──────────────────────────────────────────────────────────────────────────
\section{The paradigm, restated}
\label{sec:paradigm}

\subsection{Statement}

\normsv{}\textbf{An external streaming protocol is a transport
view over an agent's dispatch sequence.} The wire shape an
external client receives is a deterministic projection of the
typed frames the agent emits during dispatch. The adapter that
produces the projection runs as an ability rooted at the Hub,
mints one envelope per request, and emits the projection as
its body.

The agent dispatch sequence is the source of truth. The OpenAI
SSE stream is one of its views. Other views may co-exist (Anthropic
\texttt{messages} streaming, Gemini \texttt{generateContent},
future protocols) without changing the source-of-truth side.
Each view is a pure filter: a different filter, a different
projection, the same dispatch.

\subsection{Why this is not a chat-API wrapper}

\derived{} A naive design would introduce a parallel
\texttt{chat\_completions} primitive in EasyNet --- a new
state\_type, a new receipt class, a new ability namespace ---
and forward to it. This is what \emph{wrapping} looks like, and
it is structurally wrong here for the same reason RFC-006-B v0.6
rejected ``new state\_type for web\_app'': it duplicates an
ontology that already exists. EasyNet already has a notion of
``the streaming output of an agent's work'' --- the dispatch
sequence the chat ability is built on. The OpenAI surface
projects that sequence; it does not host an alternative one.

\subsection{Why a chat-base ability is already a streaming
completion}

\derived{} Every agent driver currently registered with EasyNet
(Claude Code via stream-json, Codex via SSE) emits a typed-frame
stream during dispatch. The frame types and an OpenAI streaming
completion's chunk types map onto each other:

\begin{center}
\renewcommand{\arraystretch}{1.2}
\begin{tabularx}{\linewidth}{L{0.30\linewidth} X}
\toprule
\textbf{Agent dispatch frame} & \textbf{OpenAI SSE chunk} \\
\midrule
\texttt{system}/\texttt{init} (driver handshake)        & dropped \\
\texttt{assistant.message.text} (incremental tokens)    & \texttt{choices[].delta.content} \\
\texttt{assistant.message.tool\_use}                    & \texttt{choices[].delta.tool\_calls[]} (when client requested tool surface; otherwise dropped or summarised) \\
\texttt{user.tool\_result}                              & dropped (internal dispatch loop) \\
\texttt{assistant.message.thinking}                     & optional; mapped to \texttt{choices[].delta.reasoning} when the client opts in \\
\texttt{result} (terminal)                              & \texttt{choices[].finish\_reason} + \texttt{[DONE]} sentinel \\
\bottomrule
\end{tabularx}
\end{center}

The mapping is deterministic, line-for-line, with no
state-bearing transformations. INV-3 makes this normative.

\subsection{The minimum reference instance: web-builder.chat}

\derived{} The smallest instance with all four ingredients
(an external protocol, a deterministic streaming filter, a
capability-style auth, a complete receipt trail) is exposing the
existing \texttt{web-builder.chat} ability as
\texttt{model = "web-builder"} on a single \texttt{/v1/chat/completions}
endpoint. Cursor / Continue / openai-python configured against
the EasyNet Hub URL with an EasyNet-issued API key call
\texttt{web-builder} the same way they would call
\texttt{gpt-4o}. The first end-to-end demo is one Cursor request
that lands inside web-builder's dispatch, picks up the agent's
real reply, and returns it byte-equivalent to the wire shape
Cursor expects.

% ──────────────────────────────────────────────────────────────────────────
\section{Assumptions}
\label{sec:assumptions}

\normsv{} v0.1 is written under the following assumptions. Each
is stated explicitly so downstream implementation work can be
discharged independently while this appendix's semantic claims
remain stable.

\begin{assumptionbox}{A-1: Chat-base abilities exist}
EasyNet has at least one ability of the form
\texttt{<agent>.chat} whose dispatch produces a streaming
sequence of typed frames (text, tool\_use, tool\_result,
result). The current implementation registers
\texttt{<agent>.chat} for each registered agent through
\texttt{chat\_ability::register} (RFC-001 Phase 3 et al.). v0.1
does not change that surface.
\end{assumptionbox}

\begin{assumptionbox}{A-2: Keyring abilities mint resource URAs}
The keyring family (RFC-002 Phase 3, vault-style) supports
ability invocations that mint long-lived capability tokens
addressed as \texttt{resource/api\_key.<id>} URAs. Whether the
mint primitive is named \texttt{<user>.keyring.mint\_api\_key},
\texttt{<user>.api\_key.create}, or something else is
implementation detail; what matters is that:

\begin{itemize}[leftmargin=2em,itemsep=0.1em]
    \item the minted URA is conformant to v4.1.5 (three URA segments,
    \texttt{resource/api\_key.<id>});
    \item presence of \texttt{<id>} at the HTTP boundary
    cryptographically authenticates the bearer (RFC-002 §3 vault
    rules apply);
    \item revocation is a single ability call.
\end{itemize}
\end{assumptionbox}

\begin{assumptionbox}{A-3: Receipts can be chained cross-class}
The receipt store can record both canonical (state-changing)
and operational (state-preserving) receipts and link them
through \texttt{causal\_context.Scalar}. INV-4's receipt trail
relies on this; it is the same machinery RFC-006-B v0.6 §5 uses
for receipt-chain audit walking.
\end{assumptionbox}

\begin{assumptionbox}{A-4: The Hub agent owns its adapter abilities}
\texttt{01HUB.openai.chat\_completions} is a hub-rooted ability
in the same sense \texttt{01HUB.pages.serve} is. The
\emph{ability ownership widens to (user, agent, hub, resource)}
amendment from RFC-006-B Phase 0 covers this case verbatim.
\end{assumptionbox}

% ──────────────────────────────────────────────────────────────────────────
\section{The four normative invariants}
\label{sec:invariants}

\subsection{INV-1 --- Adapter Purity}

\begin{invariant}[INV-1: Adapter Purity (OpenAI)]
\texttt{ability/01HUB.openai.chat\_completions}, the ability
that bridges the OpenAI streaming protocol into the invocation
graph, \textbf{MUST} be a pure transport adapter. Specifically:
\begin{itemize}[leftmargin=2em,itemsep=0.1em]
    \item It \textbf{MUST NOT} mutate any agent state.
    \item It \textbf{MUST} emit \emph{exactly one} canonical
    receipt per OpenAI request, against the API-key resource
    (INV-4); no other canonical receipts originate from this
    adapter.
    \item It \textbf{MAY} emit operational receipts (per-frame
    or aggregated; lossy observability).
    \item It \textbf{MUST} delegate the dispatch entirely to
    the underlying chat-base ability via the standard
    \texttt{forward\_invoke} dispatch path. It \textbf{MUST NOT}
    embed an alternative dispatcher.
    \item It \textbf{MUST NOT} cache, summarise, or transform
    the agent's reply beyond mechanical wire-shape framing
    (chunking by SSE \texttt{data:} lines; closing with
    \texttt{[DONE]}).
\end{itemize}
\end{invariant}

\derived{} Adapter Purity rules out the entire failure mode in
which the OpenAI surface accidentally becomes a parallel
authoritative state plane: a Hub implementation that adds
``output filtering'' or ``response rewriting'' beyond the
wire-shape transformation has stopped being an adapter and is
now an authority. INV-1 forbids it at the specification layer,
not at the implementation layer.

\subsection{INV-2 --- Capability-URI Key}

\begin{invariant}[INV-2: API Key as Capability Resource]
Every API key issued for OpenAI compatibility \textbf{MUST}
correspond to a v4.1.5-conformant URA of the form
\texttt{resource/api\_key.<id>}, where \texttt{<id>} is an
unguessable token of at least 256 bits of entropy. The URA
\textbf{MUST}:
\begin{itemize}[leftmargin=2em,itemsep=0.1em]
    \item be minted by a keyring ability (A-2);
    \item resolve, via the keyring, to a \emph{single}
    \texttt{user/<username>} URA --- the user the key
    represents at admission time;
    \item expose at least one revocation ability so the key
    can be deactivated without restarting the daemon;
    \item carry zero conversation state of its own (the key is
    the capability; it does not hold messages, threads, or
    cached context).
\end{itemize}
The HTTP \texttt{Authorization: Bearer <id>} header is the
transport rendering of the URA. Presence of the header is
\emph{the} authentication signal --- no OAuth, no token
exchange, no scope object.
\end{invariant}

\begin{notebox}
The minimal implementation gives each key a flat ``acts-as
\texttt{user/<username>}'' grant. Future revisions may attach
narrower scopes (e.g., ``can call only models matching
\texttt{<owner>.\*.chat}'') by recording them on the key
resource at mint time. Scope refinement is additive;
v0.1 ships flat keys.
\end{notebox}

\subsection{INV-3 --- Filter Determinism}

\begin{invariant}[INV-3: Stream Filter Determinism]
The function from agent dispatch frame sequence to OpenAI SSE
chunk sequence \textbf{MUST} be a pure deterministic projection.
For a given dispatch sequence \(D\) and a given filter
configuration \(F\):
\[
    \texttt{sse}(D, F) = \texttt{filter}(D, F)
\]
where \texttt{filter} contains no time, no random, no host-
specific data, no agent identity beyond what is already in
\(D\), and no model-output post-processing.

It follows that:
\begin{itemize}[leftmargin=2em,itemsep=0.1em]
    \item Two Hub implementations on different hosts running
    the same filter configuration against the same dispatch
    sequence \textbf{MUST} produce byte-identical SSE streams
    (modulo HTTP-layer fragmentation that does not change the
    bytes after re-assembly).
    \item A receipt-chain auditor walking from the canonical
    auth receipt down through dispatch receipts \textbf{MUST}
    be able to recompute the SSE bytes the client received,
    given the dispatch frames the receipts cite.
    \item The filter \textbf{MUST NOT} introduce dynamic content
    (timestamps, IP-derived data, server-name banners). Any such
    content belongs in the dispatch sequence itself, not in the
    projection.
\end{itemize}
\end{invariant}

\derived{} INV-3 is what makes OpenAI compat federate. Without
it, two Hubs claiming the same API key resource may serve
different bytes for the same request; with it, a CDN, a logging
proxy, or an auditor can rely on the dispatch sequence as the
canonical record and treat the SSE bytes as a function of it.

\begin{notebox}
INV-3 does not forbid all dynamic content in chat replies; it
forbids dynamic content from \emph{the filter}. The agent's own
dispatch may legitimately include timestamps the LLM produced
(``it is now 2026-05-04''). What INV-3 forbids is the adapter
adding such content unilaterally.
\end{notebox}

\subsection{INV-4 --- Auth Receipt Trail}

\begin{invariant}[INV-4: Auth Receipt Trail]
Every OpenAI request received at \texttt{/v1/chat/completions}
\textbf{MUST}:
\begin{enumerate}[leftmargin=2em,itemsep=0.1em]
    \item Mint one canonical receipt with:
    \begin{itemize}[leftmargin=1.5em,itemsep=0.05em]
        \item \texttt{ability\_uri = ability/01HUB.openai.chat\_completions}
        \item \texttt{caller\_uri = hub/01HUB}
        \item \texttt{callee\_uri = hub/01HUB}
        \item \texttt{subject\_uri = resource/api\_key.<id>}
        \item \texttt{state\_transition.body\_hash = sha256(canonicalised
        request body)}
        \item \texttt{state\_transition.usage = \{ prompt\_tokens,
        completion\_tokens, total\_tokens, model \}}
    \end{itemize}

    \item Forward dispatch through the resolved chat-base
    ability with \texttt{causal\_context.Scalar = <id of
    canonical receipt above>} so each dispatch-side receipt
    is chain-walkable back to the auth receipt.

    \item Honour an \texttt{api\_key} revocation that lands
    after admission but before dispatch terminates by surfacing
    a structured stream-end event (the OpenAI client receives a
    \texttt{[DONE]} after a final \texttt{finish\_reason =
    "cancelled"} chunk).
\end{enumerate}
\end{invariant}

\derived{} INV-4 is the adapter's narrow but real
authoritative-state contribution: ``a request was authorised
under this key at this moment with this body hash and these
usage numbers.'' Everything else stays in dispatch land.

\subsection{Why these four, no fewer, no more}

\derived{} Following the degeneracy method of AXIOM.tex
§why-six and RFC-006-B v0.6 §invariants:

\begin{itemize}[leftmargin=2em,itemsep=0.18em]
    \item \textbf{Without INV-1}, the OpenAI surface is a second
    authoritative state plane. An operator with access to the
    Hub process can rewrite chat replies without going through
    the dispatch path; the receipt chain is no longer the single
    source of truth.
    \item \textbf{Without INV-2}, API keys carry their own
    semantics outside the URA system. Revocation, scoping,
    auditing all become bespoke. Federation across realms
    becomes impossible because a key issued by one realm has
    no canonical name in another.
    \item \textbf{Without INV-3}, the same model on two hosts
    returns different bytes for the same request. CDN caching,
    third-party audit, and replay verification all collapse.
    \item \textbf{Without INV-4}, an OpenAI request enters
    EasyNet without a primary ledger entry. The receipt chain
    has gaps: ``the agent dispatched, but on whose authority?''
\end{itemize}

Each invariant blocks a class of failure that the others do not
address; removing any one collapses the model.

% ──────────────────────────────────────────────────────────────────────────
\section{State transition model in detail}
\label{sec:transitions}

\subsection{Mint key: \texttt{<user>.keyring.mint\_api\_key}}

\normsv{} The user-rooted keyring ability that mints an OpenAI
API key for that user. Canonical, state-changing.

\begin{lstlisting}
caller   = caller_agent_uri
callee   = caller_agent_uri                   (self-call)
ability  = ability/<user>.keyring.mint_api_key
subject  = user/<user>                         (the user the key acts as)
args     = {
    label:       "cursor on alice's mac",     (optional human label)
    expires_at:  null | "<rfc3339>"           (null = no auto-expiry)
}
\end{lstlisting}

\textbf{Pre-state}: zero or more \texttt{api\_key} resources owned by the user.

\textbf{Post-state}: one new \texttt{resource/api\_key.<id>} resource exists, owned by the user, status = \texttt{active}.

\textbf{Receipt class}: canonical. The receipt's
\texttt{state\_transition.post\_state} carries the minted URA;
\texttt{state\_transition.token\_hash} carries
\(\text{sha256}(\text{token\_bytes})\) so the chain proves a
specific bearer existed without storing the secret.

\subsection{Authorise: \texttt{01HUB.openai.chat\_completions}}

\normsv{} The Hub-rooted adapter ability invoked by the HTTP
boundary on every request. Canonical (one auth receipt per
call) but adapter-pure (no other state mutation).

\begin{lstlisting}
caller   = hub/01HUB                           (Hub-as-caller)
callee   = hub/01HUB                           (self-call)
ability  = ability/01HUB.openai.chat_completions
subject  = resource/api_key.<id>               (parsed from Authorization)
args     = {
    request: {
        model:    "<ability URA tail>",        (e.g. "web-builder")
        messages: [...],
        stream:   true,
        ...
    }
}
\end{lstlisting}

\textbf{Pre-state}: \texttt{resource/api\_key.<id>} exists with status = \texttt{active}.

\textbf{Post-state}: same. No state mutation in the adapter.

\textbf{Receipt class}: canonical (per INV-4). The receipt
records body hash and usage; subsequent dispatch receipts cite
this receipt as \texttt{causal\_context.Scalar}.

\subsection{Dispatch: \texttt{<agent>.chat}}

\normsv{} The chat-base ability the OpenAI request resolved to.
Operational, state-preserving (the agent's reply is a function
of input + agent state at this instant; no canonical state is
mutated by reading the reply).

\begin{lstlisting}
caller   = hub/01HUB
callee   = agent/<resolved-from-model-name>
ability  = ability/<agent>.chat
subject  = user/<resolved-from-api-key>
args     = { prompt: <last user message>, context: <other messages> }
nonce    = <derived from the canonical auth receipt id>
causal_context = Scalar(<canonical auth receipt id>)
\end{lstlisting}

The dispatch's own typed-frame stream is what the filter projects to OpenAI SSE.

\textbf{Receipt class}: operational; one per call, with frame-level operational sub-receipts as the dispatch produces tool calls / sub-invocations.

\subsection{Revoke: \texttt{<user>.keyring.revoke}}

\normsv{} Revocation primitive. Canonical, state-changing.

\begin{lstlisting}
caller   = caller_agent_uri
callee   = caller_agent_uri
ability  = ability/<user>.keyring.revoke
subject  = resource/api_key.<id>
args     = { reason: "user-initiated" | "compromised" | ... }
\end{lstlisting}

\textbf{Pre-state}: \texttt{api\_key} resource exists with status = \texttt{active}.

\textbf{Post-state}: same resource exists with status = \texttt{revoked}, \texttt{revoked\_at} populated. Subsequent OpenAI requests with that bearer are rejected at the auth step (HTTP 401) with one canonical \emph{rejection} receipt for audit.

\section{Receipt semantics}
\label{sec:receipts}

\normsv{} v0.1 keeps the two-tier classification of RFC-006-B
v0.6 and adds the precise inter-receipt linkage INV-4 mandates:

\begin{center}
\renewcommand{\arraystretch}{1.18}
\begin{tabularx}{\linewidth}{l X X}
\toprule
\textbf{Tier} & \textbf{For transitions that are} & \textbf{Used by} \\
\midrule
Canonical
 & \texttt{mint\_api\_key}, \texttt{revoke}, \texttt{openai.chat\_completions} (auth entry per request)
 & receipt-chain audit; \texttt{prev\_receipt\_hash} chain; persistent \\
Operational
 & \texttt{<agent>.chat} dispatch; per-frame observability
 & lossy ring buffer; aggregable rates surfaced through \texttt{observe.health} \\
\bottomrule
\end{tabularx}
\end{center}

\subsection{Cross-tier linkage}

A canonical OpenAI auth receipt and the operational dispatch
receipts it triggered are linked by:

\begin{itemize}[leftmargin=2em,itemsep=0.12em]
    \item dispatch receipt's \texttt{causal\_context.Scalar} = canonical auth receipt's invocation\_id;
    \item canonical auth receipt's \texttt{state\_transition.dispatch\_root\_invocation\_id} = the root operational invocation id (not the receipt id; the invocation id pins the chain head).
\end{itemize}

A receipt walker traversing from any operational dispatch
receipt can therefore reach the canonical auth receipt in one
hop (\texttt{causal\_context}) and recover the API-key URA, the
body hash, and the usage numbers.

\subsection{What receipts witness, in this reference instance}

\normsv{}

\begin{itemize}[leftmargin=2em,itemsep=0.18em]
    \item \textbf{mint\_api\_key receipt} witnesses ``user \(U\) created
    API key resource \(R\) at time \(t\), token hash \(H\),
    label \(L\), expiry \(E\)''.
    \item \textbf{auth receipt (per request)} witnesses ``a
    request bearing key \(R\) for model \(M\) with body hash
    \(H_b\) and usage \(\{i, o, t\}\) entered the system at
    \(t'\), forwarded into dispatch root invocation
    \(I_{root}\)''.
    \item \textbf{dispatch receipts (per call)} witness the
    agent's tool calls + sub-invocations; each cites the auth
    receipt; together they reconstruct what \texttt{<agent>.chat}
    actually did under this key.
    \item \textbf{revoke receipt} witnesses ``key \(R\) revoked
    at \(t''\), reason \(\rho\)''.
\end{itemize}

A consumer can therefore answer ``how many requests did key
\(R\) make and what did they cost?'' from canonical receipts
alone, and ``what tools did the agent invoke under each request?''
by walking down into operational receipts.

% ──────────────────────────────────────────────────────────────────────────
\section{URL = invocation, in this reference instance}
\label{sec:url}

\subsection{Endpoint surface}

\normsv{} The Hub binds three endpoints under the realm's HTTPS
listener:

\begin{lstlisting}
GET   /v1/models
POST  /v1/chat/completions
POST  /v1/embeddings                    (deferred; see §10)
\end{lstlisting}

Only \texttt{/v1/chat/completions} is in scope for v0.1 normative
behaviour. \texttt{/v1/models} is included because the
OpenAI-shape clients enumerate it for model discovery; v0.1
returns the realm's chat-base ability registry projected as a
list of model names. \texttt{/v1/embeddings} is reserved for a
future v0.2 once an embedding-base ability primitive lands.

\subsection{The transformation}

\normsv{} An HTTP request:

\begin{lstlisting}
POST /v1/chat/completions HTTP/1.1
Authorization: Bearer easynet-sk-<id>
Content-Type: application/json
{
  "model":    "web-builder",
  "messages": [{"role":"user","content":"hi"}],
  "stream":   true
}
\end{lstlisting}

is translated by \texttt{ability/01HUB.openai.chat\_completions} into:

\begin{lstlisting}
caller    = hub/01HUB
callee    = hub/01HUB                       (auth ability)
ability   = ability/01HUB.openai.chat_completions
subject   = resource/api_key.<id>           (parsed from Bearer)
args      = { model: "web-builder", messages: [...], ... }
\end{lstlisting}

which then forwards (per INV-1) into:

\begin{lstlisting}
caller    = hub/01HUB
callee    = agent/web-builder
ability   = ability/web-builder.chat
subject   = user/<resolved-from-api-key>
args      = { prompt: "hi", context: ... }
causal_context = Scalar(<canonical auth receipt id>)
\end{lstlisting}

The dispatch's frames are filtered (per INV-3) into SSE chunks
streamed back as the HTTP response body.

\subsection{Why this is not a routing rule}

\derived{} Per INV-1, the Hub's adapter does not own the
dispatch; it forwards. Per INV-3, the bytes the adapter writes
to the SSE stream are a deterministic projection of the
dispatch frames. Together, INV-1 + INV-3 mean the OpenAI URL
form is a projection of the URA form: every OpenAI request has
a URA, every OpenAI response equals the URA's invocation result
projected through \texttt{filter}, and the correspondence is
verifiable through the receipt chain.

This is the load-bearing claim of the paradigm extension. The
OpenAI API is no longer ``a competing protocol'' or ``an
external dependency''; it is a viewing surface onto the
invocation graph, the same way HTTP file fetch is.

% ──────────────────────────────────────────────────────────────────────────
\section{Reference implementation (informative)}
\label{sec:ref-impl}

\textit{This section is non-normative. It witnesses that the
v0.1 specification is realisable on current EasyNet-Cli
infrastructure. Conforming implementations are not required to
follow these specific choices.}

\subsection{New abilities (Hub + keyring)}

Three abilities to register:

\begin{lstlisting}
src/runtime/agents/openai_compat/
  mod.rs              ← register `01HUB.openai.chat_completions` +
                        `01HUB.openai.list_models`
  filter.rs           ← stream-json -> OpenAI SSE projection
                        (deterministic; pure fn)
  auth.rs             ← parse Authorization header; resolve
                        API-key URA via keyring; emit canonical
                        auth receipt
  models.rs           ← /v1/models — project the local ability
                        registry's chat-base entries

src/runtime/keyring/
  api_key_ability.rs  ← <user>.keyring.mint_api_key /
                        <user>.keyring.revoke /
                        <user>.keyring.list_api_keys
  api_key_storage.rs  ← persist to ~/.easynet/keyring/api_keys.toml
                        (reuse RFC-002 vault primitives)

src/runtime/hub/
  openai_listener.rs  ← extend the Hub axum router with
                        /v1/{chat/completions,models}; route to the
                        Hub abilities above
\end{lstlisting}

\subsection{The filter, sketched}

\normsv{} The filter consumes stream-json events and emits OpenAI
SSE chunks. Pseudocode:

\begin{lstlisting}
async fn filter(stream: StreamJsonRx, sink: SseSink) {
    let mut completion_id = format!("chatcmpl-{}", random_token());
    let model = current_model_name();
    let created = current_request_time_pin();   // pinned at request
                                                // start; constant
                                                // for the whole stream
                                                // so INV-3 holds.

    sink.send_first_chunk(completion_id, model, created, "assistant");

    while let Some(frame) = stream.recv().await {
        match frame {
            FrameType::AssistantText(text) => {
                sink.send_delta_chunk(completion_id, model, created, text);
            }
            FrameType::AssistantToolUse(tool) if client_supports_tools => {
                sink.send_tool_call_chunk(completion_id, model, created, tool);
            }
            FrameType::AssistantThinking(t) if client_opted_in_reasoning => {
                sink.send_reasoning_chunk(completion_id, model, created, t);
            }
            FrameType::Result { stop_reason } => {
                sink.send_finish_chunk(completion_id, model, created, stop_reason);
                sink.send_done_sentinel();
                return;
            }
            _ => { /* drop: handshake, tool_result, etc. */ }
        }
    }
}
\end{lstlisting}

\subsection{The MVP demo}

\normsv{} The minimum existence proof is one Cursor request that
hits the Hub, authenticates against an EasyNet-issued API key,
forwards to \texttt{web-builder.chat}, and returns the agent's
real reply with byte-equivalent OpenAI wire framing. This
discharges all four invariants simultaneously: Adapter Purity
(the Hub did not invent the reply), Capability-URI Key (the
Bearer matched a \texttt{resource/api\_key.<id>} URA),
Filter Determinism (a second Cursor request with identical body
returned identical SSE), Auth Receipt Trail (the receipt store
shows one canonical auth receipt per request, dispatch
receipts citing it).

\subsection{Test stratification}

\normsv{} Three matrices, each layered three deep, mirroring
RFC-006-B v0.6 §6:

\begin{itemize}[leftmargin=2em,itemsep=0.15em]
    \item \textbf{Matrix A (unit/integ, in-process)}: filter
    determinism (same input frames $\to$ same SSE bytes), key
    minting + revocation transitions, malformed-Authorization
    rejection paths.
    \item \textbf{Matrix B (CLI surface, out-of-process)}:
    \texttt{easynet api-key \{create,list,revoke\}} subcommands.
    \item \textbf{Matrix C (Docker e2e, real container, real
    OpenAI client)}: cursor + openai-python smoke runs against
    a containerised Hub; SSE round-trip; tool-use rendering;
    revocation mid-stream.
\end{itemize}

% ──────────────────────────────────────────────────────────────────────────
\section{Beyond MVP --- post-v0.1 work, in shape only}
\label{sec:post-mvp}

\deferred{}

\begin{itemize}[leftmargin=2em,itemsep=0.15em]
    \item \textbf{Embeddings} (\texttt{POST /v1/embeddings}).
    Requires an embedding-base ability primitive
    (\texttt{<agent>.embed} or similar). Same paradigm: the
    response is a deterministic projection of the agent's
    embedding output, the API key auth is identical.
    \item \textbf{Anthropic-shape adapter}
    (\texttt{POST /v1/messages}). A second filter, same
    dispatch source. Models the same paradigm against a
    second external protocol.
    \item \textbf{Gemini-shape adapter}
    (\texttt{POST /v1beta/models/<model>:streamGenerateContent}).
    A third filter.
    \item \textbf{Scoped API keys}: keys whose admission scope
    is narrower than ``acts as user $U$''. E.g., key may invoke
    only models in a specific ability namespace, or only at
    rates below a threshold. Scope is stored on the key
    resource at mint time; admission consults it.
    \item \textbf{Usage / billing receipts}: enrich the auth
    receipt's \texttt{state\_transition.usage} with cost
    attribution per ability invocation, so an audit can
    answer ``what did this user cost the cluster across all
    chat-base abilities last month?''.
    \item \textbf{Tool-call surface}: explicit OpenAI-shape
    tool registration so a client can pass \texttt{tools=[...]}
    and have the agent's tool\_use frames map onto the client's
    tool descriptors. v0.1 either drops tool\_use frames or
    summarises them as text; the proper mapping is a v0.2 item.
    \item \textbf{Conversation continuation across requests}.
    The OpenAI protocol is stateless (every request carries the
    full \texttt{messages} history); EasyNet's chat ability is
    too. A client that wants ``server-side history'' must
    introduce a separate resource (e.g.
    \texttt{resource/conversation.<id>}) and reference it
    explicitly --- RFC-006-D's territory.
\end{itemize}

Each is purely additive. None changes the four invariants of
v0.1; they instantiate the paradigm in new ways or layer
additional resources on top.

% ──────────────────────────────────────────────────────────────────────────
\section{Open decisions}
\label{sec:open}

\begin{decisionbox}
\textbf{API-key bearer prefix.} v0.1 specifies
\texttt{Bearer easynet-sk-<id>} as the user-visible token shape,
with \texttt{<id>} a 256-bit unguessable hex/base32 token. The
\texttt{easynet-sk-} prefix is convention only --- it lets a
secret-scanner (e.g., GitHub's leaked-credential detector)
recognise EasyNet keys distinctly from OpenAI's
\texttt{sk-proj-} or Anthropic's \texttt{sk-ant-}. The exact
prefix string is non-normative.
\end{decisionbox}

\begin{decisionbox}
\textbf{Model name resolution.} v0.1 resolves
\texttt{model = "web-builder"} to ability
\texttt{web-builder.chat} by suffix-appending \texttt{.chat}.
A future revision may permit explicit URAs (\texttt{model =
"easynet:///r/easynet.run/ability/web-builder.chat"}) for
disambiguation across realms, or short aliases stored on the
hub. Both fit the paradigm; v0.1 picks the simplest form.
\end{decisionbox}

\begin{decisionbox}
\textbf{Tool-use rendering.} Mapping agent tool\_use frames to
OpenAI \texttt{tool\_calls} requires a tool descriptor schema
the client agreed to up-front (via the \texttt{tools=[]}
request field). v0.1's normative position is: \emph{drop}
tool\_use frames from the SSE projection unless the request
explicitly listed compatible tools, and \emph{always} record
them in the operational receipt chain. v0.2 specifies the
descriptor handshake.
\end{decisionbox}

\begin{decisionbox}
\textbf{Revocation propagation latency.} INV-2 requires
revocation to take effect ``without restarting the daemon'';
INV-4 requires mid-stream revocation to surface as a
\texttt{[DONE]} after \texttt{finish\_reason = "cancelled"}.
The implementation needs a watch on the keyring's revocation
stream; the latency budget is post-MVP (v0.2).
\end{decisionbox}

% ──────────────────────────────────────────────────────────────────────────
\section{Coverage check}

\begin{center}
\renewcommand{\arraystretch}{1.18}
\begin{tabularx}{\linewidth}{l X l}
\toprule
\textbf{Rule} & \textbf{Discharged by} & \textbf{Where} \\
\midrule
URA v4.1.5 §A.URA-7 (3-segment URA)        & \texttt{resource/api\_key.<id>} parses to three URA segments  & \S\ref{sec:invariants} \\
URA v4.1.5 §9 (callee \(\in\) hub/device/agent) & callee on auth = \texttt{hub/01HUB}; callee on dispatch = \texttt{agent/<resolved>}     & \S\ref{sec:url} \\
ONTOLOGY §3.2.1 (ability owner \(\in\) user/agent/hub/resource) & \texttt{01HUB.openai.chat\_completions} (hub-rooted), \texttt{<user>.keyring.mint\_api\_key} (user-rooted)         & \S\ref{sec:transitions} \\
AXIOM (every invocation produces a receipt) & two-tier classification + cross-tier linkage                       & \S\ref{sec:receipts} \\
RFC-006-B v0.6 INV-1 (Adapter Purity)       & inherited verbatim as RFC-006-C INV-1                               & \S\ref{sec:invariants} \\
RFC-006-B v0.6 INV-3 (Deterministic Projection) & inherited verbatim as RFC-006-C INV-3                            & \S\ref{sec:invariants} \\
RFC-002 keyring vault (capability identity) & API keys minted by keyring abilities; addressed as resource URAs    & \S\ref{sec:invariants} \\
forward\_invoke endgame contract            & adapter forwards via standard \texttt{forward\_invoke}              & \S\ref{sec:url} \\
\bottomrule
\end{tabularx}
\end{center}

\paragraph{Closing.} Every chat-base ability is, after a
deterministic stream filter, an OpenAI streaming completion.
Every API key is, by construction, a v4.1.5-conformant
\texttt{resource/api\_key.<id>} URA capability. The OpenAI
endpoint is one ability call from a Hub-rooted adapter into a
chat-base agent ability, and the wire bytes the client receives
are a verifiable projection of the dispatch frames the agent
emits. EasyNet does not host LLM endpoints; it lets cursor /
continue / openai-python view the invocation graph through the
glass cursor / continue / openai-python already speak.

\vfill
\hrule
\vspace{0.4em}
\noindent\small
\textit{The OpenAI protocol is not a competitor and not a
dependency. It is a viewing surface --- one of many --- onto
agents whose work, with or without the surface, is recorded
in EasyNet's invocation graph.}

\end{document}
